\chaptertitle{The Evolution of Cooperation}

                           The Prisoner's Dilemma
                         Computer Tournaments and
                         the Evolution of Cooperation
                                         May, 1983

LIFE is filled with paradoxes and dilemmas. Sometimes it even feels as if the essence of
living is the sensing-indeed, the savoring-of paradox. Although all paradoxes seem
somehow related, some paradoxes seem abstract and philosophical, while others touch on
life very directly. A very lifelike paradox is the so-called "Prisoner's Dilemma",
discovered in 1950 by Melvin Dresher and Merrill Flood of the RAND Corporation.
Albert W. Tucker wrote the first article on it, and in that article he gave it its now-famous
name. I shall here present the Prisoner's Dilemma-first as a metaphor, then as a formal
problem.
The original formulation in terms of prisoners is a little less clear to the uninitiated, in my
experience, than the following one. Assume you possess copious quantities of some item
(money, for example), and wish to obtain some amount of another item (perhaps stamps,
groceries, diamonds). You arrange a mutually agreeable trade with the only dealer of that
item known to you. You are both satisfied with the amounts you will be giving and
getting. For some reason, though, your trade must take place in secret. Each of you agrees
to leave a bag at a designated place in the forest, and to pick up the other's bag at the
other's designated place. Suppose it is clear to both of you that the two of you will never
meet or have further dealings with each other again.
Clearly, there is something for each of you to fear: namely, that the other one will leave
an empty bag. Obviously, if you both leave full bags, you will both be satisfied; but
equally obviously, getting something for clothing is even more satisfying. So you are
tempted to leave an empty bag. In fact, you can even reason it through quite rigorously
this way: "If the dealer brings a full bag, I'll be better off having left an empty bag,
because I'll have gotten




The Prisoner's Dilemma Computer Tournaments and the Evolution of Cooperation              715

all that I wanted and given away nothing. If the dealer brings an empty bag, I'll be better
off having left an empty bag, because I'll not have been cheated I'll have gained nothing
but lost nothing either. Thus it seems that no matte what the dealer chooses to do, I'm
better off leaving an empty bag. So I'll lease an empty bag."
The dealer, meanwhile, being in more or less the same boat (though at the other end of
it), thinks analogous thoughts and comes to the parallel conclusion that it is best to leave
an empty bag. And so both of you, with your impeccable (or impeccable-seeming( logic,
leave empty bags, and go away empty-handed. How sad, for if you had both just
cooperated, you could have each gained something you wanted to have. Does logic
prevent cooperation? This is the issue of the Prisoner's Dilemma.

                                         *   *   *

In case you're wondering why it is called "Prisoner's Dilemma", here's the reason.
Imagine that you and an accomplice (someone you have no feelings for one way or the
other) committed a crime, and now you've both been apprehended and thrown in jail, and
are fearfully awaiting trials. You are being held in separate cells with no way to
communicate. The prosecutor offers each of you the following deal (and informs you
both that the identical deal is being offered to each of you-and that you both know that as
well'): "We hake a lot of circumstantial evidence on you both. So if you both claim
innocence, we will convict you anyway and you'll both get two years in jail. But if you
will help us out by admitting your guilt and making it easier for us to convict your
accomplice-oh, pardon me, your alleged accomplice why, then, we'll let you out free.
And don't worry- about revenge-your accomplice will be in for five years! How about
it?" Warily you ask, "But what if we both say we're guilty?" "Ah, well, my, friend-I'm
afraid you'll both get four-year sentences, then."
        Now you're in a pickle! Clearly, you don't want to claim innocence if your partner
has sung, for then you're in for five long years. Better you should both have sung-then
you'll only get four. On the other hand, if your partner claims innocence, then the best
possible thing for you to do is sing. since then you're out scot-free! So at first sight, it
seems obvious what you should do: Sing! But what is obvious to you is equally obvious
to your opposite number, so now it looks like you both ought to sing, which means -Sing
Sing for four years! At least that's what logic tells you to do. Funny\% since if both of you
had just been illogical and maintained innocence, you d both be in for only half as long!
Ah, logic does it again.

                                         *   *   *

Let us now go back to the original metaphor and slightly alter its conditions. Suppose that
both you and your partner very much want to have




The Prisoner's Dilemma Computer Tournaments and the Evolution of Cooperation            716

a regular supply of what the other has to offer, and so, before conducting your first
exchange, you agree to carry on a lifelong exchange, once a month. You still expect
never to meet face to face. In fact, neither of you has any idea how old the other one is, so
you can't be very sure of how long this lifelong agreement may go on, but it seems safe to
assume it'll go on for a few months anyway, and very likely for years.
         Now, what do you do on your first exchange? Taking an empty bag seems fairly
nasty as the opening of a relationship-hardly an effective way to build up trust. So
suppose you take a full bag. and the dealer brings one as well. Bliss-for a month. Then
you both must go back. Empty, or full? Each month, you have to decide whether to defect
(take an empty bag) or to cooperate (take a full one). Suppose that one month,
unexpectedly, your dealer defects. Now what do you do? Will you suddenly decide that
the dealer can never be trusted again, and from now on always bring empty bags, in effect
totally giving up on the whole project forever? Or will you pretend you didn't notice, and
continue being friendly? Or-will you try to punish the dealer by some number of
defections of your own? One? Two? A random number? An increasing number,
depending on how many defections you have experienced? Just how mad will you get?
         This is the so-called iterated Prisoner's Dilemma. It is a very difficult problem. It
can be, and has been, rendered more quantitative and in that form studied with the
methods of game theory and computer simulation. How does one quantify it? One builds
a payoff matrix presenting point values for the various alternatives. A typical one is
shown in Figure 29-1a. In this matrix, mutual cooperation earns both parties 2 points (the
subjective value of receiving a full hag of what you need while giving up a full hag of
what you have). 'Mutual defection earns you both 0 points (the subjective value of
gaining nothing and losing nothing, aside from making a vain trip out to the forest that
month). Cooperating while the other defects stings: you get -1 point while the rat gets 4
points! Why so many? Because it is so pleasurable to get something for nothing. And of
course, should you happen to be a rat some month when the dealer has cooperated, then
you get 4 points and the dealer loses 1.
         It is obvious that in a collective sense, it would be best for both of you to always
cooperate. But suppose you have no regard whatsoever for the other. There is no
''collective good" you are both working for. You are both supreme egoists. Then what?
The meaning of this term, "egoist", can perhaps be made clear by the following. Suppose
you and your dealer have developed a trusting relationship of Mutual cooperation over
the years, when one day you receive secret and reliable information that the dealer is
quite sick and will soon die, probably within a month or two. The dealer has no reason to
suspect that you have heard this. Aren't you highly tempted to defect, all of a sudden,
despite all your years of cooperating? You are, after all, out for yourself and no one else
in this cruel, cruel world. And since it seems that this may very well be the dealer’s last
month, why not profit




The Prisoner's Dilemma Computer Tournaments and the Evolution of Cooperation              717

        FIGURE 29-1. The Prisoner's Dilemma.

        In (a), a Prisoner's Dilemma payoff matrix in the case of a dealer and a buyer of
commodities or services, in which both participants have a choice: to cooperate (i.e., to deliver
the goods or the payment) or to defect (i.e., to deliver nothing). The numbers attempt to represent
the degree of satisfaction of each partner in the transaction.
        In (b), the formulation of the Prisoner's Dilemma to which it owes its name: in terms of
prisoners and their opportunities for double-crossing or collusion. The numbers are negative
because they represent punishments: the length of both prisoners' prospective jail sentences, in
years. This metaphor is due to Albert W Tucker.
        In (c), a Prisoner's Dilemma formulation where all payoffs are nonnegative numbers.
This is my canonical version, following the usage in Robert Axelrod's book, The Evolution of
Cooperation.

as much as possible from your secret knowledge? Your defection may never be punished,
and at the worst, it will be punished by one last-gasp defection by the dying dealer.
        The surer you are that this next turn is to be the very last one, the more you feel
you must defect. Either of you would feel that way, of course, on learning that the other
one was nearing the end of the rope. This is what is meant by “egoism”. It means you
have no feeling of friendliness or




The Prisoner's Dilemma Computer Tournaments and the Evolution of Cooperation                  718

goodwill or compassion for the other player; you have no conscience; all you care about
is amassing points, more and more and more of them.
        What does the payoff matrix for the other metaphor, the one involving prisoners,
look like? It is shown in Figure 29-1b. The equivalence of this matrix to the previous
matrix is clear if you add a constant-namely, 4-to all terms in this one. Indeed, we could
add any constant to either matrix and the dilemma would remain essentially unchanged.
So let us add 5 to this one so as to get rid of all negative payoffs. We get the canonical
Prisoner's Dilemma payoff matrix, shown in Figure 29-1c. The number 3 is called the
reward for mutual cooperation, or R for short. The number 1 is called the punishment, or
P. The number 5 is T, the temptation, and 0 is S, the sucker's payoff. The two conditions
that make a matrix represent a Prisoner's Dilemma situation are these:

                                      (I) T>R>P>S

                                     (2) (T+S)/2 < R

        The first one simply makes the argument go through for each of you, that "it is
better for me to defect no matter what my counterpart does". The second one simply
guarantees that if you two somehow get locked into out-of-phase alternations (that is,
"you cooperate, I defect" one month and "you defect, I cooperate" the next), you will not
do better-in fact, you will do worse-than if you were cooperating with each other each
month.
        Well, what would be your best strategy? It can be shown quite easily that there is
no universal answer to this question. That is, there is no strategy that is better than all
other strategies under all circumstances. For consider the case where the other player is
playing ALL D-the strategy of defecting each round. In that case, the best you can
possibly do is to defect each time as well, including the first. On the other hand, suppose
the other player is using the Massive Retaliatory Strike strategy, which means "I'll
cooperate until you defect and thereafter I'll defect forever." Now if you defect on the
very first move, then you'll get one T and all P's thereafter until one of you dies. But if
you had waited to defect, you could have benefited from a relationship of mutual
cooperation, amassing many R's beforehand. Clearly that bunch of R's will add up to
more than the single T if the game goes on for more than a few moves. This means that
against the ALL D strategy, ALL D is the best counterstrategy, whereas "Always
cooperate unless you learn that you or the other player is just about to die, in which case
defect" is the best counterstrategy against Massive Retaliatory Strike. This simple
argument shows that how you should play depends on who you're playing.
        The whole concept of the "quality" of a strategy takes on a decidedly more
operational and empirical meaning if one imagines an ocean populated by dozens of little
beings swimming around and playing Prisoner's Dilemma over and over with each other.
Suppose that each time two such beings encounter each other, they recognize each other
and




The Prisoner's Dilemma Computer Tournaments and the Evolution of Cooperation           719

remember how previous encounters have gone. This enables each one to decide what it
wishes to do this tine. Now if each organism is continually swimming around and
bumping into the others, eventually, each one will have met even other one numerous
times, and thus all strategies will have been given the opportunity to interact with each
other. By "interact", what is meant here is certainly not that anyone knocks anyone else
out of the ocean, as in an elimination tournament. The idea is simply that each organism
gains zero or more points in each meeting, and if sufficient time is allowed to elapse,
everybody will have met with everybody else about the same number of times. and now
the only question is: Who has amassed the most points? Amassing points is truly the
name of the game.
         It doesn't matter if you have ''beaten'' anyone, in the sense of haying gained more
Front interacting with them than they gained from interacting with you. That kind of
"victory" is totally irrelevant here. What matters is not the number of "victories" rung tip
by any individual. but the individual's total point count-a number that measures the
individual's overall viability in this particular "sea" of many strategies. It sounds nearly
paradoxical. but you could lose many-indeed, all-of your individual skirmishes, and yet
still come out the overall winner.
         As the image suggests very strongly, this whole situation is highly relevant to
questions in evolutionary biology. Can totally selfish and unconscious organisms living
in a common environment come to evolve reliable cooperative strategies Can cooperation
emerge in a world of pure egoists? In a nutshell, can cooperation evolve out of
noncooperation? If so, this has revolutionary import for the theory of evolution, for many
of its critics have claimed that this was one place that it was hopelessly snagged.
         Well, as it happens, it has now been demonstrated rigorously and definitively that
such cooperation can emerge, and it was done through a computer tournament conducted
by political scientist Robert Axelrod of the Political Science Department and the Institute
for Public Policy Studies of the University of Michigan in Ann Arbor. More accurately,
Axelrod first studied the ways that cooperation evolved by means of a computer
tournament, and when general trends emerged, he was able to spot the underlying
principles and prove theorems that established the facts and conditions of cooperation's
rise from nowhere. Axelrod has written a fascinating and remarkably thought-provoking
book on his findings, called The Evolution of Cooperation, published in 1984 by Basic
Books, Inc. (Quoted sections below are taken front an early draft of that book.)
Furthermore, he and evolutionary biologist William 1), Hamilton have worked out and
published many of the implications of these discoveries for evolutionary theory. Their
work has won much acclaim-including the 198 1 Newcomb Cleveland prize, a prize
awarded annually by the American Association for the Advancement of Science for “an
outstanding paper published in Science”.




The Prisoner's Dilemma Computer Tournaments and the Evolution of Cooperation            720

There are really three aspects of the question "Can cooperation emerge in a world of
egoists?'' The first is: How can it get started at all? the second is: Can cooperative
strategies survive better than their non-cooperative rivals? The third one is: Which
cooperative strategies will do the best, and how will they come to predominate?
         To make these issues vivid, let me describe Axelrod's tournament and its
somewhat astonishing results. In 1979, Axelrod sent out invitations to a number of
professional game theorists, including people who had published articles on the Prisoner's
Dilemma. telling them that he wished to pit many strategies against one another in a
round-robin Prisoner's Dilemma tournament, with the overall goal being to amass as
many points as possible. He asked for strategies to be encoded as computer programs that
could respond to the 'C' or 'D' of another player, taking into account the remembered
history of previous interactions with that same player. A program should always reply
with a 'C' or a 'I)', of course, but its choice need not be deterministic. That is. consultation
of a random-number generator was allowed at any point in a strategy.
         Fourteen entries were submitted to Axelrod, and he introduced into the field one
more program called RANDOM, which in effect flipped a coin (computationally
simulated, to be sure) each move, cooperating if heads came up, defecting otherwise. The
field was a rather variegated one, consisting of programs ranging from as few as four
lines to as many as 77 lines (of Basic). Every program was made to engage each other
program (and a clone of itself) 200 times. No program was penalized for running slowly.
The tournament was actually run five times in a row, so that ,pseudo-effects caused by
statistical fluctuations in the random-number generator would be smoothed out by
averaging.
         The program that won was submitted by the old Prisoner's Dilemma hand, Anatol
Rapoport, a psychologist and philosopher from the University of Toronto. His was the
shortest of all submitted programs, and is called TIT FOR TAT'.. TIT FOR TAT uses a
very simple tactic:

       Cooperate on move 1;
       Thereafter, do whatever the other player did the previous move.

         That is all. It sounds outrageously simple. How in the world could such a program
defeat the complex stratagems devised by other experts?
         Well, Axelrod claims that the game theorists in general did not go far enough in
their analysis. They looked "only two levels deep", when in fact they should have looked
three levels deep to do better. What precisely does this mean? He takes a specific case to
illustrate his point. Consider the entry called JOSS (submitted by Johann Joss, a
mathematician from Zurich, Switzerland). JOSS´s strategy is very similar to TIT FOR
TAT´s, in that it




The Prisoner's Dilemma Computer Tournaments and the Evolution of Cooperation               721

begins by cooperating, always responds to defection by defecting and nearly always
responds to cooperation by cooperating. The hitch is that JOSS uses a random-number
generator to help it decide when to pull a "surprise defection" on the other player. JOSS
is set up so that it has a 10 percent probability of defecting right after the other player has
cooperated.
         In playing TIT FOR TAT, JOSS will do fine until it tries to catch TIT FOR TIT
off guard. When it defects, TIT FOR TAT retaliates with a single defection, while JOSS
"innocently" goes hack to cooperating. Thus we have a "DC" pair. On the next move, the
'C' and 'D' will switch places since each program in essence echoes the other's latest
move, and so it will go: CD, then DC, CD, DC, and so on. There may ensue a long
reverberation set off by JOSS's D, but sooner or later, JOSS will randomly, throw in
another unexpected D after a C from TIT FOR TAT. At this point, there will he a "DD"
pair, and that determines the entire rest of the match. Both will defect forever, now. The
"echo" effect resulting from JOSS's first attempt at exploitation and TIT FOR TAT's
simple punitive act lead ultimately to complete distrust and lack of cooperation.
         This may seem to imply that both strategies are at fault and will suffer for it at the
hands of others, but in fact the one that suffers from it most is JOSS, since JOSS tries out
the same trick on partner after partner, and in mans cases this leads to the same type of
breakdown of trust, whereas TIT FOR TA T, never defecting first, will never be the
initial cause of a breakdown of trust. Axelrod's technical term for a strategy that never
defects before its opponent does is nice. TIT FOR TAT is a nice strategy,JOSS is not.
Note that "nice" does not mean that a strategy never defects! TIT FOR TAT defects when
provoked, but that is still considered being "nice".
         Axelrod summarizes the first tournament this way:

         A major lesson of this tournament is the importance of minimizing echo effects in an
environment of mutual power. A sophisticated analysis must go at least three levels deep. First is
the direct effect of a choice. This is easy, since a defection always earns more than a cooperation.
Second are the indirect effects, taking into account that the other side may or may not punish a
defection. This much was certainly appreciated by many of the entrants. But third is the fact that
in responding to the defections of the other side, one may be repeating or even amplifying one's
own previous exploitative choice. Thus a single defection may be successful when analyzed for
its direct effects, and perhaps even when its secondary effects are taken into account. But the real
costs may be in the tertiary effects when one's own isolated defections turn into unending mutual
recriminations. Without their realizing it, many of these rules actually wound up punishing
themselves. With the other player serving as a mechanism to delay the self-punishment by a few
moves, this aspect of self-punishment was not perceived by the decision rules ....
         The analysis of the tournament results indicates that there is a lot to be learned about
coping in an environment of mutual power. Even expert strategists from political science,
sociology, economics, psychology, and mathematics made the systematic errors of being too
competitive for their own




The Prisoner's Dilemma Computer Tournaments and the Evolution of Cooperation                   722

good. not forgiving enough, and too pessimistic about the responsiveness of the other side.

        Axelrod not only analyzed the first tournament, he even performed a number of
"subjunctive replays" of it, that is, replays with different sets of entries. He found, for
instance, that the strategy called TIT FOR TWO T.ATS, which tolerates two defections
before getting mad (but still only- strikes back once), would have won, had it been in the
line-up. Likewise, two other strategies he discovered, one called REVISED DOWNING
and one called LOOK-AHEAD, would have come in first had they been in the
tournament.
        In summary, the lesson of the first tournament seems to have been that it is
important to be nice ("don't be the first to defect") and forgiving, ("don't hold a grudge
once you've vented your anger"). TIT FOR TAT possesses both these qualities, quite
obviously.

                                            *   *   *

        After this careful analysis, Axelrod felt that significant lessons had been
unearthed, and he felt convinced that more sophisticated strategies could be concocted,
based on the new information. Therefore he decided too hold a second, larger computer
tournament. For this tournament, he not only invited all the participants in the first round,
but also advertised in computer hobbyist magazines, hoping to attract people who were
addicted to programming and who would be willing to devote a good deal of time to
working out and perfecting their strategies. To each person who entered, Axelrod sent a
full and detailed analysis of the first tournament, along with a discussion of the
"subjunctive replays" and the strategies that would have won. He described the strategic
concepts of "niceness" and "forgiveness" that seemed to capture the lessons of the
tournament so well, as well as strategic pitfalls to avoid. Naturally, each entrant realized
that all the other entrants had received the same mailing, so that everyone knew that
everyone knew that everyone knew that ...
        There was a large response to Axelrod's call for entries. Entries were received
from six countries, from people of all ages, and from eight different academic disciplines.
Anatol Rapoport entered again, resubmitting TIT FOR TAT (and was the only one to do
so, even though it was explicitly stated that anyone could enter any program written by
anybody). A ten-year-old entered, as did one of the world's experts on game theory and
evolution, John Maynard Smith, professor of biology at the University of Sussex in
England, who submitted TIT FOR TWO TATS. Two people separately submitted
REVISED DOWNING.
        Altogether, 62 entries were received, and generally speaking, they were of a
considerably higher degree of sophistication than those in the first tournament. The
shortest was again TIT FOR TAT, and the longest was a program from New Zealand,
consisting of 152 lines of Fortran. Once again,




The Prisoner's Dilemma Computer Tournaments and the Evolution of Cooperation                  723

RANDOM was added to the field, and with a flourish and a final carriage return, the
horses were off' Several hours of computer time later, the results came in.


        The outcome was nothing short of stunning: TIT FOR 7:I T, the simplest program
submitted, won again. What's more, the two programs submitted that had won the
subjunctive replays of the first tournament now turned up way down in the list: TIT FOR
TWO TATS came in 24th, and REVISED DOWNING ended up buried in the bottom
half of the field.
        This may seem horribly nonintuitive, but remember that a program's success
depends entirely on the environment in which it is swimming. There is no single "best
strategy" for all environments, so that winning in one tournament is no guarantee of
success in another. TIT FOR TAT has the advantage of being able to "get along well"
with a great variety of strategies, while other programs are more limited in their ability to
evoke cooperation. Axelrod puts it this way:

        What seems to have happened is an interesting interaction between people who drew one
lesson and people who drew another lesson from the first round. Lesson One was "Be nice and
forgiving." Lesson Two was more exploitative: "If others are going to be nice and forgiving, it
pays to try to take advantage of them." The people who drew Lesson One suffered in the second
round from those who drew Lesson Two.

        Thus the majority of participants in the second tournament really had not grasped
the central lesson of the first tournament: the importance of being willing to initiate and
reciprocate cooperation. Axelrod feels so strongly about this that he is reluctant to call
two strategies playing against each other "opponents"; in his book he always uses neutral
terms such as "strategies" or "players". He even does not like saying they are playing
against each other, preferring "with". In this article, I have tried to follow his usage, with
occasional departures. One very striking fact about the second tournament is the success
of "nice" rules: of the top fifteen finishers. only one (placing eighth) was not nice.
Amusingly, a sort of mirror image held: of the bottom fifteen finishers, only one was
nice!
        Several non-nice strategies featured rather tricky probes of the opponent (sorry!),
sounding it out to see how much it "minded" being defected against. Although this kind
of probing by a program might fool occasional opponents, more often than not it
backfired, causing severe breakdowns of trust. Altogether, it turned out to be very costly
to try to use defections to "flush out" the other player's weak spots. It turned out to be
more profitable to have a policy of cooperation as often as possible, together with a
willingness to retaliate swiftly against any attempted undercutting. Note, however, that
strategies featuring massive retaliation were less successful than TIT FOR TAT with its
more gentle policy of restrained retaliation.




The Prisoner's Dilemma Computer Tournaments and the Evolution of Cooperation              724

Forgiveness is the key here, for it helps to restore the proverbial "atmosphere of mutual
cooperation" (to use the phrase of international diplomacy) after a small skirmish.
        "Be nice and forgiving" was in essence the overall lesson of the first tournament.
Apparently, though, many people just couldn't get themselves to believe it, and were
convinced that with cleverer trickery and scheming, they could win the day. It took the
second tournament to prove them dead wrong. And out of the second tournament, a third
key strategic concept emerged: that of provocability-the notion that one should "get mad"
quickly at defectors. and retaliate. Thus a more general lesson is: "Be nice, provocable,
and forgiving."
                                           *       *       *
        Strategies that do well in a wide variety of environments are called by Axelrod
robust. and it seems that ones with "good personality traits"-that is, nice. provocable. and
forgiving strategies-are sure to be robust. TIT FOR TAT is by no means the only possible
strategy with these traits, but it is the canonical example of such a strategy, and it is
astonishingly robust.
        Perhaps the most vivid demonstrations of TIT FOR TAT's robustness were
provided by various subjunctive replays of the second tournament. The principle behind
any replay involving a different environment is quite simple. From the actual playing of
the tournament, you have a 63X63 matrix documenting how well each program did
against each other program. Now, the effective "population" of a program in the
environment can be manipulated mathematically by attaching a weight factor to all that
program's interactions, then just retotaling all the columns. This way you can get
subjunctive instant replays without having to rerun the tournament.
        This simple observation means that the results of a huge number of potential
subjunctive tournaments are concealed in, but potentially extractable from, the 63X63
matrix of program-z s.-program totals. For instance, Axelrod discovered, using statistical
analysis, that there were essentially six classes of strategies in the second tournament. For
each of these classes, he conducted a subjunctive instant replay of the, tournament by
quintupling the importance (the weight factor) of that class alone, thus artificially
inflating a certain strategic style's population in the environment. When the scores were
retotaled, TIT FOR TAT emerged victorious in five out of six of those hypothetical
tournaments, and in the sixth it placed second.
        Undoubtedly the most significant and ingenious type of subjunctive replay that
Axelrod tried was the ecological tournament. Such a tournament consists not merely of a
single subjunctive replay, but of -a whole cascade of hypothetical replays, each one's
environment determined by the results of the previous one. In particular, if you take a
program's score in a tournament as a measure of its "fitness", and if you interpret
"fitness" as meaning "number of progeny in the next generation", and finally, if you let




The Prisoner's Dilemma Computer Tournaments and the Evolution of Cooperation             725

"next generation" mean "next tournament", then what you get is that each tournament's
results determine the environment of the next one-arid in particular, successful programs
become more copious in the next tournament. This type of iterated tournament is called
"ecological" because it simulates ecological adaptation (the shifting of a fixed set of
species' populations according to their mutually defined and dynamically developing
environment), as contrasted with evolution via mutation (where new species can come
into existence).
        As one carries an ecological tournament through generation after generation, the
environment gradually changes. In a paraphrase of how Axelrod puts it, here is what
happens. At the very beginning, poor programs and good programs alike are equally
represented. As time passes, the poorer ones begin to drop out while the good ones
flourish. But the tank order of the good ones may now change. because their "goodness"
is no longer being measured against the same field of competitors as initially. Thus
success breeds ever more success-but only provided that the success derives from
interaction with other similarly successful programs. If, by contrast, some program's
success is due mostly to its ability to milk "dumber" programs for all they're worth, then
as those programs are gradually squeezed out of the picture, the exploiter's base of
support will be eroded and it will suffer a similar fate.
        A concrete example of ecological extinction is provided by HARRINGTON. the
only non-nice program among the top fifteen finishers in the second tournament. In the
first 200 generations of the ecological tournament, while TIT FOR TAT and other
successful nice programs were gradually increasing their percentage of the population,
HARRINGTON ' too was increasing its percentage. This was a direct result of
HARRINGTON´S exploitative strategy. However, by the 200th generation, things began
to take a noticeable turn. Weaker programs were beginning to go extinct, which meant
fewer and fewer dupes for HARRINGTON to profit from. Soon the trend became
apparent: HARRINGTON could not keep up with its nice rivals. By the 1,000th
generation, HARRINGTON was as extinct as the dodos it had exploited. Axelrod
summarizes:

        Doing well with rules that do not score well themselves is eventually a self'-defeating
process. Not being nice may look promising at first, but in the long run it can destroy the very
environment it needs for its own success.

         Needless to say, TIT FOR TAT fared spectacularly well in the ecological
tournament, increasing its lead ever more. After 1,000 generations, not only was TIT
FOR TAT ahead, but its rate of growth was greater than that of any other program. This
is an almost unbelievable success story,, all the more so because of the absurd simplicity
of the "hero". One amusing aspect of it is that 7'IT FOR TAT did not defeat a single one
of its rivals in their encounters. This is not a quirk; it is in the nature of TIT FOR TAT.
TIT FOR TAT simply cannot defeat anyone; the best it can achieve is a tie, and often it
loses (though not by much).




The Prisoner's Dilemma Computer Tournaments and the Evolution of Cooperation               726

Axelrod makes this point very clear:

TIT FOR TAT won the tournament, not by beating the other placer, but by eliciting behavior
from the other player which allowed both to do well. TIT FOR TAT was so consistent at eliciting
mutually rewarding outcomes that it attained a higher overall score than any other strategy in the
tournament.
        So in a non-zero-sum world you do not have to do better than the other player to do well
for yourself. This is especially true when you are interacting with many different players. Letting
each of them do the same or a little better than you is fine, as long as you tend to do well yourself.
There is no point in being envious of the success of the other player, since in an iterated Prisoner's
Dilemma of long duration the other's success is virtually a prerequisite of your doing well for
yourself.

He gives examples from everyday life in which this principle holds. Here is one:

A firm that buys from a supplier can expect that a successful relationship will earn profit for the
supplier as well as the buyer. There is no point in being envious of the supplier's profit. Any
attempt to reduce it through an uncooperative practice, such as by not paying your bills on time,
,will only encourage the supplier to take retaliatory action. Retaliatory action could take many
forms, often without being explicitly labeled as punishment. It could be less prompt deliveries,
lower quality control, less forthcoming attitudes on volume discounts, or less timely news of
anticipated market conditions. The retaliation could make the envy quite expensive. Instead of
worrying about the relative profits of the seller, the buyer should worry about whether another
buying strategy would be better.

Like a business partner who never cheats anyone, TIT FOR TAT never beats anyone-yet
both do very well for themselves.

        One idea that is amazingly counterintuitive at first in the Prisoner's Dilemma is
that the best possible strategy to follow is ALL D if the other player is unresponsive. It
might seem that some form of random strategy might do better, but that is completely
wrong. If I have laid out all my moves in advance, then playing TIT FOR TAT will do
you no good, nor will flipping a coin. You should simply defect every move. It matters
not what pattern I have chosen. Only if I can be influenced by your play will it ever do
you any good to cooperate.
        Fortunately, in an environment where there are programs that cooperate (and
whose cooperation is based on reciprocity), being unresponsive is a very poor strategy,
which in turn means that ALL D is a very poor, strategy. The single unresponsive
competitor in the second tournament was RAND0m and it finished next to last. The last-
place finishers strategy was responsive, but its behavior was so inscrutable that it looked
unresponsive




The Prisoner's Dilemma Computer Tournaments and the Evolution of Cooperation                     727

And in a more recent computer tournament conducted by Marek Lugowsl, and myself in
the Computer Science Department at Indiana L-university three ALL-D's came in at the
very bottom (out of 53), with a couple of RA.VDO.11's giving them a tough fight for the
honor.
       One way to explain TIT FOR TIT's success is simply to say that it elicits
cooperation, via friendly persuasion. Axelrod spells this out as follows:

        Part of its success might be that other rules anticipate its presence and are
designed to do well with it. Doing well with TIT FOR TAT requires cooperating with it,
and this in turn helps TIT FOR TIT Even rules that were designed to see what then could
get away with quickly apologize to TIT FOR F-IT. Am rule that tries to take advantage of
TIT FOR T.AT will simple hurt itself. TIT FOR TAT benefits from its own
nonexploitability because three conditions are satisfied:

       1. The possibility of encountering TIT FOR TAT is salient;
       2. Once encountered. TIT FOR TAT is easy to recognize: and
       3. Once recognized, TIT FOR TAT 's nonexploitability is easy to appreciate

This brings out a fourth "personality trait" (in addition to niceness. provocability, and
forgiveness) that may play an important role in success: recognizability. or
straightforwardness. Axelrod chooses to call this trait clarity, and argues for it with
clarity:

       Too much complexity can appear to be total chaos. If you are using a strategy that
appears random, then you also appear unresponsive to the other player. If you are
unresponsive, then the other player has no incentive to cooperate with you. So being so
complex as to be incomprehensible is yen dangerous.

       How rife this is with morals for social and political behavior! It is rich food for
thought.

                                             * *    *

         Anatol Rapoport cautions against overstating the advantages of TIT FOR TIT; in
particular, he believes that TIT FOR LIT is too harshly retaliatory on occasion. It can also
be persuasively argued that TIT FOR TAT is too lenient on other occasions. Certainly
there is no evidence that TI T FOR TIT is the ultimate or best possible strategy. Indeed,
as has been emphasized repeatedly, the very concept of "best possible" is incoherent,
since all depends on environment. In the tournament at Indiana University mentioned
earlier, several TIT-FOR-TAT- like strategies did better than pure TIT FOR LIT did.
They all shared, however, the three critical "character traits" whose desirability had been
so clearly delineated by Axelrod's prior analysis of the important properties of TIT FOR
TAT They were simple a little better than TIT FOR Liz' at detecting nonresponsiveness,
and when they were convinced the other player was unresponsive, they switched over to
an ALL-D mode.




The Prisoner's Dilemma Computer Tournaments and the Evolution of Cooperation            728

In his book, Axelrod takes pains to spell out the answers to three fundamental questions
concerning the temporal evolution of cooperation in a world of raw egoism. The first
concerns initial viability: How can cooperation get started in a world of unconditional
defection-a "primordial sea" swarming with unresponsive .ILL-D creatures? The answer
(whose proof I omit here) is that invasion by small clusters of conditionally cooperating
organisms, even if they form a tiny minority, is enough to give cooperation a toehold.
One cooperator alone will die, but small clusters of cooperators can arrive (via mutation
or migration, say) and propagate even in a hostile environment, provided they are
defensive like TIT FOR TAT. Complete pacifists-Quaker-like programs-will not survive,
however, in this harsh environment.
         The second fundamental question concerns robustness: What type of strategy does
well in unpredictable and shifting environments? We have already seen that the answer to
this question is: Any strategy possessing the four fundamental "personality traits" of
niceness, provocability, forgiveness, and clarity. This means that such strategies, once
established. will tend to flourish, especially in an ecologically evolving world.
         The final question concerns stability: Can cooperation protect itself from
invasion? Axelrod proved that it can indeed. In fact, there is a gratifying asymmetry to
his findings: Although a world of "meanies" (beings using the inflexible ALL-D strategy)
is penetrable by cooperators in clusters, a world of cooperators is not penetrable by
meanies, even if they arrive in clusters of any size. Once cooperation has established
itself, it is permanent. As Axelrod puts it, "The gear wheels of social evolution have a
ratchet."
         The term "social" here does not mean that these results necessarily apply only to
higher animals that can think. Clearly, four-line computer programs do not think-and yet,
it is in a world of just such "organisms" that cooperation has been shown to evolve. The
only "cognitive abilities" needed by TIT FOR T.-IT are: (1) recognition of previous
partners, and (2) memory of what happened last time with this partner. Even bacteria can
do this, by interacting with only one other organism (so that recognition is automatic) and
by responding only to the most recent action of their "partner" (so that memory
requirements are minimal). The point is that the entities involved can be on the scale of
bacteria, small animals, large animals, or nations. There is no need for "reflective
rationality"; indeed, TIT FOR LIT could be called "reflexive" (iii the sense of being as
simple as a knee-jerk reflex) rather than "reflective".


        For people who think that moral behavior toward others can emerge only when
there is imposed some totally external and horrendous threat (say, of the fire-and-
brimstone sort) or soothing promise of heavenly reward (such as eternal salvation), the
results of this research must give pause for thought. In one sentence, Axelrod captures the
whole idea: Mutual cooperation can




The Prisoner's Dilemma Computer Tournaments and the Evolution of Cooperation           729

emerge in a world of egoists without central control, by starting with a cluster ol
individuals who rely on reciprocity.
        There are so man situations in the world today where these ideas seem of extreme
relevance-indeed, urgency-that it is very tempting to draw morals all over the place. In
the later chapters of his book, Axelrod offers advice about how to promote cooperation in
human affairs, and at the very end the political scientist in him cautiously ventures some
broad conclusions concerning global issues, which are a fitting way for me to conclude as
well

         Today, the most important problems facing humanity are in the arena of
 international relations where independent, egoistic nations face each other in a state of
 near anarchy. Many of these problems take the form of an iterated Prisoner's
 Dilemma. Examples can include arms races, nuclear proliferation. crisis bargaining,
 and military escalation. Of course, a realistic understanding of these problems would
 have to take into account mane factors not incorporated into the simple Prisoner's
 Dilemma formulation, such as ideology. bureaucratic politics, commitments,
 coalitions, mediation, and leadership. Nevertheless, we can use all the insights we can
 get.
         Robert Gilpin [in his book War and Change in World Politics] points out that
 from the ancient Greeks to contemporary scholarship all political theory addresses
 one fundamental question: "How can the human race. whether for selfish or more
 cosmopolitan ends, understand and control the seemingly blind forces of history?" In
 the contemporary world this question has become especially acute because of the
 development of nuclear weapons.
         The advice given in this book to players of the Prisoner's Dilemma might also
 serve as good advice to national leaders as well: Don't be envious, don't he the first to
 defect, reciprocate both cooperation and defection, and don't be too clever. Likewise,
 the techniques discussed in this book for promoting cooperation in the Prisoner's
 Dilemma might also be useful in promoting cooperation in international politics.
         The core of the problem is that trial-and-error learning is slow and painful.
 The conditions may all be favorable for long-run developments, but we may not have
 the time to wait for blind processes to move us slowly towards mutually rewarding
 strategies based upon reciprocity. Perhaps if we understand the process better, we can
 use our foresight to speed up the evolution of cooperation.

       Post Scriptum.
        In the course of writing this column and thinking the ideas through. I was forced
to confront over and over again the paradox that the Prisoner's Dilemma presents. I found
that I simply could not accept the seemingly flawless logical conclusion that says that a
rational player in a noniterated situation will always defect. In turning this over in my
mind and trying to articulate my objections clearly, I found myself inventing variation
after




The Prisoner's Dilemma Computer Tournaments and the Evolution of Cooperation           730

variation after variation on the basic situation. I would like to describe just a few here.
         A version of the dealer-and-buyer scenario involving bags exchanged in a forest
actually occurs in a more familiar context. Suppose I take my car in to get the oil
changed. I know little about auto mechanics, so when I come in to pick it up. I really
have no way to verify if they've done the job. For all I know, it's been sitting untouched
in their parking lot all day, and as I drive off they may be snickering behind my back. On
the other hand, maybe I've got the last laugh, for how do they know if that check I gave
them will bounce?
         This is a perfect example of how either of us could defect, but because the
situation is iterated, neither of us is likely to do so. On the other hand, suppose I'm on my
way across the country and have some radiator trouble near Gillette. Wyoming, and stop
in town to get my radiator repaired there. There is a decent chance now that one party or
the other will attempt to defect, because this kind of situation is not an iterated one. I'll
probably never again need the services of this garage, and they'll never get another check
from me. In the most crude sense, then, it's not in my interest to give them a good check,
nor is it in theirs to fix my car. But do I really defect? Do I give out bad checks? No. Why
not?
         Consider this related situation. Late at night, I bang into someone's car in a
deserted parking lot. It's apparent to me that nobody witnessed the incident. I have the
choice of leaving a note, telling the owner who's to blame, or scurrying off scot-free.
Which do I do? Similarly, suppose I-have given a lecture in a classroom in a university I
am visiting for one day, and have covered the board with chalk. Do I take the trouble of
erasing the board so that whoever comes in the next morning won't have to go to that
trouble? Or do I just leave it?


        I was recently waiting to board an airplane when a voice announced: "Passengers
holding seats in rows 24 to 36 may now board." Well, my seat was in row 4, so I waited.
A few minutes later, the voice said that passengers in rows 18 to 36 were free to board. A
group of people got up and went in. Then after a couple of minutes, rows 10 to 36 were
told they could board. A dozen people or so remained in the waiting area. For a while, we
were all patient, waiting for the final announcement allowing us to board, but after about
five minutes, people started fidgeting a bit and edging up toward the gate. Then, after
another two or three minutes, a couple of people just went right on. And then the rest of
us wondered, "Should we get on, too? Will we be left behind?" For most of the people,
the answer was obvious: they rushed to board. And once they had boarded, then the rest
of us felt kind of like suckers, and we just got on too. In effect, there was a stampede that
converted cooperators into defectors. Even the people who triggered the stampede had
originally been cooperating, but after a while the temptation




The Prisoner's Dilemma Computer Tournaments and the Evolution of Cooperation             731

got to be too great, and they broke down. At that point, some sort of phase transition, or
collective shift, took place, and the stable state of patient cooperation collapsed into a
chaotic scrambling for places. Actually, it wasn't that bad, and there was a good reason
for the relatively polite way we did board, defectors though we were: we all had seat
assignments, so it didn't matter who got on first. But imagine if the earliest defectors were
sure to get the best remaining seats! The contemporary aphorist .Ashleigh Brilliant has
found just the right bons mots to describe this sort of dilemma:

Should I abide by the rules until they're changed, or help speed tht change by breaking
them?

       Better start rushing before the rush begins!

In pondering the Prisoner's Dilemma, I could not help but be reminded of horrible
scenarios in Nazi concentration camps, where large herds of unarmed people would be
led to their deaths by small herds of armed people. It seems that a stampede by the
masses could quickly have overcome a small number of guards, at least in certain critical
narrow passageways here and there. The trouble is, it would require certain death on the
part of a few ultra-cooperators. in exchange for the liberation of a large number of other
people. Generally speaking, individuals are not willing to perform such an exchange.
Nobody wants to be in the front lines of a protest demonstration facing troops with
machine guns. Everyone wants to be in the rear. But not everyone can be in the rear! If
nobody is willing to be in the front lines, then there will be no front lines, and
consequently no demonstration at all.


        Driving a car has a certain primitive quality to it that brings out the animal in us
all, and probably that's why it confronts us with Prisoner's-Dilemma like situations so
often-more often than any other activity I can think of. How about those annoying drivers
who, when there's a long line at a freeway exit, zoom by all the politely lined-up cars and
then butt in at the very last moment, getting off 50 cars ahead of you? Are you angry t
such people. or do you do it too? Or, worse-do you do it and yet resent others who have
such gall?
        I have been struck by the relative savagery of the driving environment in the
Boston area. I know of no other city in which people are so willing to take the law into
their own hands, and to create complete anarchy. There seems to be less respect for such
things as red lights. stop signs, lines in the street, speed limits, other people's cars, and so
forth, than in any other city or state, or country that I have ever driven in. This incessant
"me-first" attitude seems to be a vicious, self-reinforcing circle. Since there are so many
people who do whatever they want, nobody can afford to be polite and let other




The Prisoner's Dilemma Computer Tournaments and the Evolution of Cooperation               732

people in ahead of them (say), for then they will be taken advantage of repeatedly and
will wind up losing totally. You simply must assert yourself in many situations, and that
means you must defect. Of course, just one defection does not an .ALL-D player make.
In fact, a retaliatory defection is just good old TIT-FOR-TAT playing. However, very
often in Boston driving, there is no way you can get back at a nasty driver who cuts in
front of you and then takes off screeching around the corner. That person is gone forever.
You can take out your frustrations only on the rest of the people near you, who are not to
blame for that driver. You can cut in ahead of them. Does this do any good? That is, does
it teach anybody a lesson? Obviously it will teach them only that it pays to defect. And
thus the spiral starts.
        Is there any way to put a halt to the descending spiral, the vortex towards
oblivion? Is there any point at which the people of Boston will collectively come to
realize that it has gotten so bad that they will all suddenly "flip" and begin to cooperate in
situations where they formerly would have defected? Can there be a stampede toward
cooperation, just as there can be a stampede toward defection?
        Clearly, if large numbers of people were to start driving much less aggressively
and nastily, everybody would benefit. Huge snarls would unsnarl-in fact would never
form. Traffic would flow smoothly and regularly. The shoulders-those favorite illegal
passing lanes for defectors-would be completely clear. So clear, in fact, that just think-
you and I could make sensational progress by swerving onto an empty shoulder and
passing everybody. Wheee! Isn't this fun? Aren't those other people suckers, staying in
the slow lane and glaring at us? Say, how come other people are barging in on us? This is
our lane. Oh, so that person in the yellow car wants to play dirty, eh? Okay. I'll show
them what playing dirty's really like!
        Sound familiar? Is there any solution to such terrible spirals? Sometimes I am
very pessimistic on that subject. Anatol Rapoport and I exchanged letters concerning this
matter, and he related a frightening anecdote. I quote from his letter:

           Do you know of the experiment performed by Martin Shubik, in which a
   dollar bill was auctioned off for S3.40? This was a consequence of a rule (the
   implications of which dawned on the subjects only when they were already
   hooked) specifying that while the highest bidder got the dollar, the second-highest
   bidder would also have to pay what he last bid. It thus became imperative to keep
   going, since the second-highest bidder (whoever he was at each stage) had
   progressively more to lose as the bids went up. Are Reagan and Andropov too
   stupid to see the point? ....
           I believe the "technological imperative'' is driving our species to
   extinction. Ever more horrendous weapons must be produced, simply because it is
   possible to produce them. Eventually the\, must be used, to justify the insane
   waste. It thus becomes imperative to seal off' the "logic" of the paradigm based on
   "deterrence", "balance of power", and similar metaphors-to make it in-assailable.




The Prisoner's Dilemma Computer Tournaments and the Evolution of Cooperation              733

   I don't think intelligence plans a part in the vicious cycle of the arms race. The
   rulers only think they make the decisions. If they were C players. they would not
   be where they are. If they started to play C while in office, they would he
   impeached. overthrown, or assassinated. Does this mean that D players are
   selected for? Possibly in the short run, but not on the time scale of evolution. H.
   sapiens is apparently not the last word, but for me, a homocentric, this is no
   consolation.

Pretty- sobering words from one of the leading rational thinkers of our era




The Prisoner's Dilemma Computer Tournaments and the Evolution of Cooperation         734

